\chapter{Cross-Chapter Practice Problems}
\label{app:practice}
The problems below are the point of the book: each one needs results from
\emph{several} chapters chained together, which is what real analysis looks like and
what the individual chapters cannot exercise on their own. Attempt each before reading
the solution, and note which chapters you had to visit.

\section{A Coupling Capacitor's Low-Frequency Loss}
\emph{Uses \cref{ch:rc,ch:splane}.} An audio stage is coupled through a capacitor into
a \SI{10}{\kilo\ohm} load, forming a high-pass. (a) Choose \(C\) for a corner at
\SI{100}{\hertz}. (b) How much is a \SI{50}{\hertz} hum component attenuated?

\subsection*{Solution}
(a) The corner of a first-order RC is \(f_c=1/(2\pi RC)\)
(\cref{ch:rc}), so
\[
C=\frac{1}{2\pi R f_c}=\frac{1}{2\pi(10^4)(100)}\approx\SI{159}{\nano\farad}.
\]
(b) The high-pass magnitude is \(|G|=\omega\tau/\sqrt{1+(\omega\tau)^2}\)
(\cref{sec:rc-freq}). At \SI{50}{\hertz}, \(\omega\tau=2\pi(50)(1.59\times10^{-3})=0.5\), so
\[
|G|=\frac{0.5}{\sqrt{1.25}}=0.447
\qquad\Longrightarrow\qquad
20\log_{10}(0.447)\approx-\SI{7.0}{\decibel}.
\]
One octave below the corner costs about \SI{7}{\decibel}---which is why a coupling
corner is placed well below the lowest wanted frequency, not at it.

\section{Designing a Trap to a Bandwidth Specification}
\emph{Uses \cref{ch:rlc,ch:filters}.} A series-LC trap must notch
\SI{7.1}{\mega\hertz} using \(L=\SI{2.2}{\micro\henry}\), and its half-power width
must be about \SI{50}{\kilo\hertz}. Find \(C\), the required \(Q\), the series
resistance that produces it, and the pole locations.

\subsection*{Solution}
From \(\omega_0=1/\sqrt{LC}\) (\cref{ch:rlc}), with
\(\omega_0=2\pi(7.1\times10^6)=\SI{4.46e7}{\radian\per\second}\),
\[
C=\frac{1}{\omega_0^2L}\approx\SI{228}{\pico\farad}.
\]
The bandwidth specification fixes \(Q\) through \(BW\approx f_0/Q\):
\[
Q=\frac{f_0}{BW}=\frac{7.1\times10^6}{5.0\times10^4}=142 .
\]
Then \(Q_s=\omega_0L/R\) gives the allowed resistance, using
\(\omega_0L=\SI{98.1}{\ohm}\):
\[
R=\frac{\omega_0L}{Q}=\frac{98.1}{142}\approx\SI{0.69}{\ohm}.
\]
Finally \(\alpha=R/2L=\SI{1.57e5}{\per\second}\) and \(\omega_d\approx\omega_0\)
(since \(\zeta=1/2Q\) is tiny), so
\[
s\approx-1.57\times10^{5}\pm\jj\,4.46\times10^{7}\ \si{\radian\per\second}.
\]
Note what the specification really demanded: less than an ohm of total loss,
\emph{including} the coil's own resistance. That is the practical reason narrow traps
need high-\(Q\) components---the pole must sit that close to the \(\jj\omega\) axis.

\section{Loading an Oscillator Tank}
\emph{Uses \cref{ch:rlcpar,ch:active}.} The same \(L=\SI{2.2}{\micro\henry}\) and
\(C=\SI{228}{\pico\farad}\) form the \emph{parallel} tank of an oscillator, with an
unloaded \(Q\) of \(200\). (a) What impedance does the tank present at resonance?
(b) The following stage loads it with \SI{5}{\kilo\ohm}. What are the loaded \(Q\) and
bandwidth, and why does it matter for phase noise?

\subsection*{Solution}
(a) For the parallel circuit \(Q_p=R_p/(\omega_0L)\) (\cref{ch:rlcpar}), so the
equivalent shunt loss resistance is
\[
R_p=Q_p\,\omega_0L=(200)(98.1)\approx\SI{19.6}{\kilo\ohm},
\]
and that is the (purely resistive) impedance at resonance---a \emph{maximum}, the
mirror of the series trap above.

(b) The load appears in parallel with \(R_p\) (\cref{ch:resistive}):
\[
R_{\text{tot}}=\frac{(19.6)(5.0)}{19.6+5.0}\approx\SI{3.98}{\kilo\ohm},
\qquad
Q_{\text{loaded}}=\frac{R_{\text{tot}}}{\omega_0L}=\frac{3980}{98.1}\approx41,
\]
so \(BW\approx f_0/Q\approx\SI{175}{\kilo\hertz}\). Loading cut \(Q\) by a factor of
five. Since \(Q\) measures how tightly the pole pair is pinned to the \(\jj\omega\)
axis (\cref{ch:active}), a heavily loaded tank lets the poles wander and the
oscillator's spectrum smears---worse phase noise. The fix is to couple the next stage
loosely, trading signal amplitude for spectral purity.

\section{An Antenna, a Line, and a Matching Section}
\emph{Uses \cref{ch:rlc,ch:filters,ch:lines}.} An antenna measures
\(Z=35+\jj40\ \si{\ohm}\) at the operating frequency, fed with \SI{50}{\ohm} coax.
(a) Find \(|\Gamma|\) and SWR. (b) The antenna is then shortened until it is resonant,
giving \(Z=\SI{35}{\ohm}\) purely real. Now what is the SWR? (c) What quarter-wave
section would match the resonant antenna exactly?

\subsection*{Solution}
(a) With \(\Gamma=(Z_L-Z_0)/(Z_L+Z_0)\) (\cref{ch:lines}),
\[
\Gamma=\frac{-15+\jj40}{85+\jj40},
\qquad
|\Gamma|=\frac{\sqrt{15^2+40^2}}{\sqrt{85^2+40^2}}=\frac{42.7}{93.9}=0.455,
\qquad
\mathrm{SWR}=\frac{1.455}{0.545}\approx2.7 .
\]
(b) At resonance the reactance is gone (\cref{ch:rlc}) and only the resistive
mismatch remains:
\[
|\Gamma|=\left|\frac{35-50}{35+50}\right|=0.176,
\qquad
\mathrm{SWR}=\frac{1.176}{0.824}\approx1.43 .
\]
Tuning out the reactance did most of the work---but not all, because
\(\SI{35}{\ohm}\ne\SI{50}{\ohm}\). This is exactly the two-part job a matching network
does (\cref{ch:filters}): cancel the reactance, then transform the resistance.

(c) The quarter-wave section needs the geometric mean (\cref{ch:lines}):
\[
Z_t=\sqrt{Z_0Z_L}=\sqrt{(50)(35)}\approx\SI{41.8}{\ohm}.
\]
Note that (b) and (c) are the two halves in order: reactance first, resistance second.

\section{What the Scope Is Really Showing You}
\emph{Uses \cref{ch:resistive,ch:rlc,ch:measurement}.} You probe a node whose
Th\'evenin resistance is \SI{47}{\kilo\ohm} with a \(10\!:\!1\) probe
(\SI{10}{\mega\ohm}, \SI{12}{\pico\farad} tip) using a three-inch ground lead.
(a) How much does the probe load the node at DC? (b) A fast edge shows overshoot and
ringing that the circuit should not produce. Where does it come from?

\subsection*{Solution}
(a) The probe is a load on the node's Th\'evenin equivalent
(\cref{ch:measurement}):
\[
\frac{V_{\text{meas}}}{V_{\text{true}}}=\frac{R_{\mathrm{in}}}{R_{\mathrm{in}}+R_{Th}}
=\frac{10\times10^6}{10\times10^6+47\times10^3}=99.5\%,
\]
so the DC reading is only \(0.5\%\) low---acceptable.

(b) The ringing is the measurement setup, not the circuit. Three inches of ground lead
is about \(L=\SI{60}{\nano\henry}\), which forms a series LC loop with the probe's
\SI{12}{\pico\farad} tip capacitance. By \(f_0=1/(2\pi\sqrt{LC})\)
(\cref{ch:rlc}),
\[
f_0=\frac{1}{2\pi\sqrt{(60\times10^{-9})(12\times10^{-12})}}\approx\SI{188}{\mega\hertz},
\]
and the loop has almost no resistance, so its \(Q\) is high and any fast edge
shock-excites it. Shorten the ground path---a one-inch spring tip moves the resonance
above \SI{350}{\mega\hertz}, out of the way.

\section{How Many Poles Does the Filter Need?}
\emph{Uses \cref{sec:decibels,ch:filters}.} A transmitter low-pass must attenuate the
second harmonic---at twice the cutoff frequency---by at least \SI{40}{\decibel}. What
Butterworth order is required?

\subsection*{Solution}
Far above cutoff an \(n\)-pole filter rolls off at \(-20n\ \si{\decibel}\) per decade
(\cref{ch:filters}). One octave is \(\log_{10}2\approx0.301\) decade, so the
attenuation an octave past cutoff is about
\[
20n\,(0.301)=6.02n\ \si{\decibel}.
\]
Setting \(6.02n\ge40\) gives \(n\ge6.6\), so \textbf{seven poles}. (The exact
Butterworth expression \(|G|^2=1/[1+(f/f_c)^{2n}]\) gives \(n=6.64\), the same
answer---the asymptotic estimate is reliable this far out.) Seven reactive elements is
a lot, which is why real transmitter filters often add \emph{zeros} instead: an
elliptic design places a transmission zero right at the harmonic and reaches
\SI{40}{\decibel} with far fewer parts (\cref{ch:filters}).

\section{An L-Network and Its Bandwidth Cost}
\emph{Uses \cref{ch:filters,ch:rlc}.} Match a \SI{50}{\ohm}
source to a \SI{1500}{\ohm} load with an L-network. Find the network \(Q\), both
reactances, and comment on the bandwidth.

\subsection*{Solution}
From \cref{ch:filters},
\[
Q=\sqrt{\frac{R_{\text{hi}}}{R_{\text{lo}}}-1}=\sqrt{\frac{1500}{50}-1}=\sqrt{29}\approx5.4,
\]
\[
X_{\text{series}}=Q\,R_{\text{lo}}=(5.4)(50)\approx\SI{269}{\ohm},
\qquad
X_{\text{shunt}}=\frac{R_{\text{hi}}}{Q}=\frac{1500}{5.4}\approx\SI{279}{\ohm}.
\]
The \(Q\) is not a free choice---the resistance ratio fixed it. And since a loaded
\(Q\) of \(5.4\) means a fractional bandwidth of roughly \(1/Q\approx18\%\)
(\cref{ch:rlc}), this match holds over a useful range but must be retuned across
bands. A \(\pi\) or T network would let you choose \(Q\) independently of the
resistance ratio, buying either wider bandwidth or sharper harmonic rejection as
needed.
